\subsection{An actual local obstruction at a bad prime}

\begin{statement}{A.55}{A ramified local point section need not be cuspidal at a bad prime}{Proved}
On the plus Fermat septic over $\Q_7$, there is an actual full local
point section whose geometrically pro-$3$ projection is intrinsically
ramified and is not cuspidal, even after passage to its geometric
abelianization. Thus the good-reduction ramification criterion A.28
does not extend as a purely local assertion to all bad primes.
This statement supplies no global full section over $\Q$ and does not
refute the literal global assertion in task A.9.
\end{statement}
\begin{proof}
There exists $y\in\Q_7$ with $y^7=-129$ and
$y\equiv46\pmod{49}$. To see this directly, put
\[
 F(z)=\frac{(4+7z)^7+129}{49}\in\Z[z].
\]
Its reduction modulo seven is $1+z$, with derivative one. Hensel's
lemma therefore gives a unique root $z\equiv6\pmod7$, and
$y=4+7z$ has the asserted properties. In particular
\[
 P=[2:y:1]\in U_7(\Q_7)
\]
is an interior point, and supplies a genuine full local point section.

Under the morphism to the tripod, its coordinate pair is
\[
 t=-X^7/Z^7=-128,\qquad 1-t=-Y^7/Z^7=129.
\]
Both are seven-adic units, with residues five and three. The cubes in
$\F_7^\times$ are exactly $\{1,-1\}$. Using the generator three of
$\F_7^\times$, their classes in
$\F_7^\times/(\F_7^\times)^3\simeq\F_3$ are respectively two and
one. The local Kummer class of the tripod section in its geometric
pro-$3$ abelianization consequently has reduction
\[
 (\kappa_3(t),\kappa_3(1-t))=(2,1)
 \quad\text{in }\F_3^2
\]
on these unit factors. There is no loss of information in this unit
test: cubing is an automorphism on $1+7\Z_7$, so a seven-adic unit is
a cube precisely when its residue is.

Every cusp section of the tripod has one of the following Kummer
forms, already modulo three:
\[
 (a,0)\quad\text{at }0,\qquad
 (0,b)\quad\text{at }1,\qquad
 (c,c+\kappa_3(-1))\quad\text{at }\infty.
\]
Here $\kappa_3(-1)=0$, since $-1$ is a cube. The pair $(2,1)$ lies
on none of these three lines. Equivalently, neither $t$ nor $1-t$
nor $(1-t)/t$ is a cube; the last ratio reduces to two. Thus the
tripod section is not cuspidal even in the pro-$3$ abelian quotient.
If the pro-$3$ projection of the section on $U_7$ were cuspidal, its
image on the tripod would have one of these cusp forms. The same
argument excludes cuspidality of its geometric abelianization.

Nevertheless this section is intrinsically ramified at seven.
Inertia acts nontrivially on $\mu_7$, hence on the twenty-one
geometric coordinate cusps of $X_7$. On the boundary submodule
\[
 \left(\bigoplus_{d\in X_7\setminus U_7}
       \Z_3(1)\right)\big/\Z_3(1)
\]
of geometric abelianization, the cyclotomic action of $I_7$ on
$\Z_3(1)$ is trivial but its permutation action is nontrivial.
As in A.48, a nonidentity permutation cannot disappear modulo the
single diagonal relation. Inner conjugation is trivial on
abelianization, so this inertia action persists for every representative
of the point section. It is ramification in the intrinsic sense of
Hoshi's Definition A.1, precisely the definition audited in A.48.

The executed program \texttt{a\_main\_bad\_prime.py} checks the
integer Hensel polynomial, six successive lifts, all finite Kummer
classes, and the three cusp lines. This example explains the actual
bad-prime limitation in A.30: a local ramification argument alone
cannot supply its claimed cusp conclusion there. The global theorem
A.31 uses the stipulated local theorem to obtain global cuspidality;
the present local point section is not a counterexample to that
distinct assertion.
\end{proof}
